\documentclass[11pt]{article}
\usepackage[margin=1in]{geometry}
\usepackage{amsmath,amssymb,amsthm,mathtools,bm,booktabs,longtable,array}
\usepackage{xcolor}
\usepackage[hidelinks]{hyperref}
\usepackage{microtype}

\hypersetup{
  pdftitle={A Four-Dimensional SU(2) Quotient-Fiber Block Theorem},
  pdfauthor={Research dossier},
  pdfsubject={Finite-block coercivity and local heat-kernel stability},
  pdfkeywords={Yang--Mills, SU(2), heat kernel, gauge quotient, renormalization group}
}

\definecolor{provedgreen}{RGB}{20,105,55}
\definecolor{conditionalblue}{RGB}{25,65,145}
\definecolor{openred}{RGB}{150,25,25}
\newcommand{\proved}{\textcolor{provedgreen}{\textsc{proved}}}
\newcommand{\conditional}{\textcolor{conditionalblue}{\textsc{conditional}}}
\newcommand{\open}{\textcolor{openred}{\textsc{open}}}
\newcommand{\R}{\mathbb R}
\newcommand{\E}{\mathbb E}
\newcommand{\Cov}{\operatorname{Cov}}
\newcommand{\rank}{\operatorname{rank}}
\newcommand{\im}{\operatorname{im}}
\newcommand{\spec}{\operatorname{spec}}
\newcommand{\su}{\mathfrak{su}(2)}
\newcommand{\W}{W_2}
\newcommand{\ind}{\mathbf 1}

\newtheorem{definition}{Definition}[section]
\newtheorem{theorem}[definition]{Theorem}
\newtheorem{proposition}[definition]{Proposition}
\newtheorem{lemma}[definition]{Lemma}
\newtheorem{corollary}[definition]{Corollary}
\theoremstyle{remark}
\newtheorem{remark}[definition]{Remark}
\newtheorem{obligation}[definition]{Open obligation}

\title{A Four-Dimensional $SU(2)$ Quotient--Fiber Block Theorem\\
\large From Gauge-Zero-Mode Removal to Local Heat-Kernel Transport}
\author{Research dossier---finite theorem, nonlinear local stability, and frontier audit}
\date{20 July 2026}

\begin{document}
\maketitle

\begin{abstract}
We construct an explicit four-dimensional $2^4$ cubical block for the tangent
theory of $SU(2)$ lattice gauge fields and separate infinitesimal gauge modes,
retained coarse curvature data, and eliminated conditional-fiber variables.
The scalar block has $81$ vertices, $216$ links, and $216$ plaquettes.  After
removing gauge directions its physical tangent space has dimension $136$; the
$24$ displayed coarse boundary-face fluxes have exact rank $17$, leaving a
$119$-dimensional scalar fiber.  Tensoring with $\su$ gives dimensions $408$,
$51$, and $357$.  A product-Hodge calculation proves
$\|d_1A\|^2\ge2\|A\|^2$ on the gauge-orthogonal quotient and hence on the
fiber.  Gaussian conditioning changes only the transported mean; the centered
fiber covariance is exactly independent of retained data.  We then prove a
finite-dimensional nonlinear stability theorem: explicit Hessian and
centered cross-gradient bounds for the local $SU(2)$ heat-kernel action imply
strong convexity and a Wasserstein fiber coefficient below one.  Heat-kernel
asymptotics and BCH analyticity ensure such bounds on a sufficiently small
local tube, but no explicit cutoff-uniform tube, global orbit atlas, defect
polymer estimate, multiscale iteration, or infrared crossover is proved.
\end{abstract}

\begin{center}
\fbox{\parbox{0.92\textwidth}{\textbf{Claim level.}
The cubical, Hodge, rank, and Gaussian results are \proved\ restricted
finite-block theorems.  The nonlinear transport implication is \proved\ from
explicit derivative bounds; its local heat-kernel realization has existential
finite-block constants.  Full-measure RG iteration, crossover, continuum
existence, and the Clay mass gap remain \open.}}
\end{center}

\tableofcontents

\section{Why this is the next target}

The exact Clay problem asks for a nontrivial axiomatic pure Yang--Mills theory
on $\R^4$ with a positive Hamiltonian gap for every compact simple gauge group
\cite{JaffeWitten}.  Earlier stages of the present programme established two
diagnostic facts.  First, translated concentrated heat kernels have
total-variation distance tending to one, so a worst-case microscopic TV basin
cannot be cutoff uniform.  Second, a block variable can respond to boundary
data because it transports a retained coarse mode even when its centered
fluctuation is independent of that data.  The local object to control is thus
the conditional fiber after gauge and retained modes have been removed.

This paper moves that idea from a two-variable abelian calibration to an
explicit four-dimensional cubical gauge complex with $\su$ color.  It does not
attempt the ultraviolet-to-infrared crossover.  The target is the finite local
theorem that must precede such an iteration.

We use the heat-kernel plaquette action because it is an exact class-function
regulator with transparent weak-field asymptotics and a long-standing role as
an alternative lattice action \cite{MenottiOnofri}.  Background minimizers and
gauge-covariant expansions are central in Balaban's RG analysis
\cite{BalabanBackground,BalabanPropagators}; our local centering is designed to
be compatible with that architecture.

\section{The frozen cubical block}

Let $I_2$ be the cellular interval with vertices $\{0,1,2\}$ and two positively
oriented edges.  The block is
\[
 B=I_2^4.
\]
Let $C^q(B;\R)$ be the real cellular $q$-cochains with the standard Euclidean
inner product.  The first two coboundaries are
\[
 d_0:C^0(B;\R)\longrightarrow C^1(B;\R),
 \qquad
 d_1:C^1(B;\R)\longrightarrow C^2(B;\R),
\]
with $d_1d_0=0$.

\begin{proposition}[Cell counts and exactness, \proved]\label{prop:counts}
The block has
\[
 |B_0|=81,\qquad |B_1|=216,\qquad |B_2|=216.
\]
Moreover,
\[
 \rank d_0=80,
 \qquad \rank d_1=136,
 \qquad \ker d_1=\im d_0.
\]
\end{proposition}
\begin{proof}
Product counting gives
\[
 |B_0|=3^4,
 \quad |B_1|=4\cdot2\cdot3^3,
 \quad |B_2|=\binom42 2^2 3^2.
\]
The one-skeleton is connected, so $\ker d_0$ consists of constants and
$\rank d_0=81-1=80$.  The cubical block is contractible, hence
$H^1(B;\R)=0$ and $\ker d_1=\im d_0$.  Rank--nullity then gives
$\rank d_1=216-80=136$.
\end{proof}

The infinitesimal link field for $SU(2)$ is
\[
 \mathcal A=C^1(B;\R)\otimes\su.
\]
The infinitesimal gauge directions are $\im d_0\otimes\su$.  We represent the
quotient by the orthogonal Coulomb slice
\begin{equation}\label{eq:H}
 \mathcal H=\ker d_0^*\otimes\su.
\end{equation}
This is a local tangent representation, not a global gauge fixing.

\begin{corollary}[Physical tangent dimension, \proved]
The scalar quotient $\ker d_0^*$ has dimension $136$.  The $\su$-valued
physical tangent space $\mathcal H$ has dimension $408$, while the gauge
subspace has dimension $240$.
\end{corollary}

\section{Retained coarse curvature and the conditional fiber}

For every coordinate pair $\mu<\nu$ and every choice of the remaining two
coordinates in $\{0,2\}$, sum $d_1A$ over the four fine plaquettes in that
$2\times2$ boundary face.  This defines
\[
 C:C^1(B;\R)\longrightarrow\R^{24}.
\]
It is invariant under linearized gauge transformations because $Cd_0=0$.

\begin{theorem}[Exact retained rank, \proved]\label{thm:rank}
The $24$ displayed coarse face fluxes have rank $17$.  Their restriction to
$\ker d_0^*$ also has rank $17$.
\end{theorem}
\begin{proof}
Let $P:C^1(I_2^4)\to C^1(I_1^4)$ sum each pair of fine boundary links into a
coarse link.  Discrete Stokes gives
$C=d_{1,\mathrm{coarse}}P$.  The map $P$ is onto.  The single coarse
hypercube has $16$ vertices and $32$ links and is contractible, so
\[
 \rank d_{1,\mathrm{coarse}}=32-(16-1)=17.
\]
Therefore $\rank C=17$.  Since $C$ vanishes on $\im d_0$ and
$C^1=\im d_0\oplus\ker d_0^*$, its restriction to the orthogonal quotient has
the same rank.  The accompanying exact rational elimination independently
certifies the rank.
\end{proof}

Choose an orthonormal identification of the image of $C|_{\ker d_0^*}$ with
$\R^{17}$ and denote the resulting independent retained map by
\[
 \overline C:\ker d_0^*\to\R^{17}.
\]
Define the scalar fiber and the $\su$ fiber by
\begin{equation}\label{eq:fiber}
 F=\ker d_0^*\cap\ker C,
 \qquad \mathcal F=F\otimes\su.
\end{equation}

\begin{corollary}[Base and fiber dimensions, \proved]
The scalar base and fiber dimensions are $17$ and $119$.  Their $\su$
counterparts have dimensions
\[
 \dim\mathcal Y=51,
 \qquad \dim\mathcal F=357.
\]
\end{corollary}

The seven dependencies among the $24$ displayed fluxes are structural coarse
Bianchi/boundary relations.  Treating all $24$ as independent would make the
conditional problem ill-typed.

\section{Certified quotient-fiber coercivity}

The tangent Yang--Mills quadratic form is
\[
 Q_t(A)=\frac1{2t}\|d_1A\|^2,
 \qquad A\in\ker d_0^*.
\]

\begin{theorem}[Product-Hodge coercivity, \proved]\label{thm:hodge}
For the block $I_2^4$,
\begin{equation}\label{eq:coercivity}
 \|d_1A\|^2\ge2\|A\|^2
 \qquad(A\in\ker d_0^*).
\end{equation}
Consequently the same inequality holds for $A\in F$ and componentwise on
$\mathcal F$.
\end{theorem}
\begin{proof}
On the interval $I_2$, the nonzero spectra of the vertex and edge Laplacians
are
\[
 \spec(d_0^*d_0)=\{0,1,3\},
 \qquad \spec(d_0d_0^*)=\{1,3\}.
\]
The Hodge Laplacian on product one-cochains is the tensor sum of one edge
Laplacian and three vertex Laplacians.  Its eigenvalue-one sector uses the
edge eigenvalue $1$ and three constant vertex modes.  These modes are exact:
they are gradients of the corresponding vertex eigenfunctions.  Hence they
are orthogonal to $\ker d_0^*$.  The next possible tensor-sum eigenvalue is
$1+1=2$.  On $\ker d_0^*$,
\[
 \Delta_1=d_0d_0^*+d_1^*d_1=d_1^*d_1,
\]
which proves \eqref{eq:coercivity}.  Restriction to $F$ can only increase the
Rayleigh minimum, and tensoring with $\su$ repeats each scalar eigenvalue three
times.
\end{proof}

The executable spectral calculation gives
\[
 \lambda_{\min}(d_1^*d_1|_{\ker d_0^*})=2,
 \qquad
 \lambda_{\min}(d_1^*d_1|_F)\approx3.438447187.
\]
The theorem uses the analytic lower bound $2$; the sharper number is a
floating-point calibration.

\section{Exact Gaussian conditional-fiber theorem}

Let $L=d_1^*d_1|_{\ker d_0^*}$ and choose independent base coordinates so
that $B:\ker d_0^*\to\R^{17}$ has full row rank.  Consider the centered
Gaussian tangent measure with precision $L/t$, conditioned on $BA=y$.

\begin{theorem}[Transported mean and invariant centered fiber, \proved]
\label{thm:gaussian}
Let
\[
 G=BL^{-1}B^*,
 \qquad M=L^{-1}B^*G^{-1}.
\]
Then
\begin{align}
 \E[A\mid BA=y]&=My,\label{eq:mean}\\
 \operatorname{Cov}(A\mid BA=y)
 &=t\bigl[L^{-1}-L^{-1}B^*G^{-1}BL^{-1}\bigr].\label{eq:cov}
\end{align}
In particular, the centered conditional variable
$Z=A-My$ has a law independent of $y$ and is supported on $F$.  Its centered
base-response coefficient is exactly zero.
\end{theorem}
\begin{proof}
The matrix $G$ is positive definite because $B$ has full row rank and $L>0$.
The standard constrained-Gaussian completion of the square gives
\eqref{eq:mean}--\eqref{eq:cov}.  Neither the covariance nor its support
depends on $y$.  Subtracting $My$ therefore produces the same fiber law for
every retained value.
\end{proof}

The computed Euclidean operator norm of $M$ is approximately $0.393174006$.
This number depends on the normalization of base coordinates and is not a
physical invariant.  The exact zero after centering is the invariant content.

\section{Local nonlinear \texorpdfstring{$SU(2)$}{SU(2)} heat-kernel stability}

Let $p_t$ be the $SU(2)$ heat kernel and let $U_e=\exp A_e$ in a sufficiently
small logarithmic chart.  The local block density, including the logarithmic
Haar Jacobian $J(A)$, has action
\begin{equation}\label{eq:action}
 S_t(A)=-\sum_{p\subset B}\log p_t(U_p(A))-\log J(A).
\end{equation}
On a local Coulomb slice, parallel transport all coarse boundary-face
holonomies to a fixed root and take logarithms.  At the trivial connection the
differential of this displayed retained map is $C\otimes I_{\su}$ and has rank
$51$.  Choose $51$ fixed independent linear combinations whose differential
is surjective.  The submersion theorem then gives local coordinates
\[
 A=\Phi(y,z),qquad y\in\R^{51},\quad z\in\R^{357},
\]
in which $y$ is retained and $z$ parameterizes the conditional fiber.

\begin{lemma}[Finite-block heat-kernel expansion, \proved\ (local)]\label{lem:hk}
There exist $t_0,\rho_0,C_B,c_B>0$ such that, for
$0<t\le t_0$ and $\|\Phi(y,z)\|_\infty\le\rho\le\rho_0$,
\begin{equation}\label{eq:hessian-remainder}
 \left\|t\nabla_z^2S_t(\Phi(y,z))-L_F\otimes I_{\su}\right\|
 \le C_B\bigl(\rho+t+e^{-c_B/t}\bigr),
\end{equation}
after a fixed bounded normalization of the local coordinates.  Here
$L_F=d_1^*d_1|_F$ at the origin.
\end{lemma}
\begin{proof}
On a compact logarithmic ball strictly inside the cut locus, the heat-kernel
parametrix and all of its required derivatives have the form
\[
 -\log p_t(\exp X)=\frac{\|X\|^2}{2t}+a_0(X)+O(e^{-c/t}),
\]
after reducing $c$ if derivatives contribute polynomial powers of $t^{-1}$.
The functions $a_0$ and the logarithmic Haar Jacobian have bounded derivatives.
The plaquette BCH expansion is
\[
 \log U_p(A)=d_1A|_p+O(\|A\|^2),
\]
with bounded derivatives on the fixed finite block.  Differentiating twice
shows that the difference from $t^{-1}d_1^*d_1$ is
$O(\rho/t)+O(1)+O(e^{-c/t}/t)$.  Multiplication by $t$ yields
\eqref{eq:hessian-remainder}.  The local coordinate map and its inverse have
bounded derivatives after shrinking $\rho_0$.
\end{proof}

\begin{theorem}[Nonlinear fiber coercivity, \proved]\label{thm:nonlinear}
Set
\[
 \eta=C_B\bigl(\rho+t+e^{-c_B/t}\bigr).
\]
Whenever $\eta<2$, the local nonlinear fiber potential is strongly convex
with
\begin{equation}\label{eq:alpha}
 \nabla_z^2S_t(\Phi(y,z))\ge\frac{2-\eta}{t}I.
\end{equation}
It consequently has a unique conditional minimizer in every convex local
fiber domain on which the bound holds.
\end{theorem}
\begin{proof}
Theorem~\ref{thm:hodge} gives
$L_F\otimes I_{\su}\ge2I$.  Apply the operator-norm perturbation bound
\eqref{eq:hessian-remainder}.  Strong convexity gives uniqueness.
\end{proof}

This is a genuine local nonlinear statement, but the constants are
existential finite-block constants.  A multiscale theorem requires explicit
values in an effective-action norm.

\section{Centered transport from derivative certificates}

Let $z_t^\star(y)$ be the conditional minimizer and recenter the common convex
fiber chart $D\subset\R^{357}$ by $w=z-z_t^\star(y)$.  Write the centered
potential as $V_y(w)$ and the normalized measure as
\[
 \widetilde\nu_y(dw)=Z_y^{-1}e^{-V_y(w)}\ind_D(w)\,dw.
\]

\begin{lemma}[Strong-convexity transport, \proved]\label{lem:transport}
Suppose $\nabla^2V_y\ge\alpha I$ on the same convex domain $D$ for all $y$, and
\[
 \sup_{w\in D}\|\nabla V_y(w)-\nabla V_{y'}(w)\|
 \le L\,d_Y(y,y').
\]
Then
\begin{equation}\label{eq:w2}
 W_2(\widetilde\nu_y,\widetilde\nu_{y'})
 \le\frac{L}{\alpha}d_Y(y,y').
\end{equation}
\end{lemma}
\begin{proof}
Strong convexity on a convex domain gives the logarithmic Sobolev inequality
with constant $1/\alpha$.  Applied to the density ratio, it bounds relative
entropy by $1/(2\alpha)$ times relative Fisher information.  The gradient
hypothesis bounds that Fisher information by $L^2d_Y(y,y')^2$.  The
logarithmic Sobolev inequality implies the quadratic transport inequality
\cite{OttoVillani}, yielding \eqref{eq:w2}.
\end{proof}

For the quadratic tangent action, centering makes $L=0$.

\begin{lemma}[Centered cross-gradient cancellation, \proved\ (local)]
\label{lem:cross}
The local coordinates may be chosen so that their tangent splitting at the
origin is the Gaussian conditional-mean direction plus the fiber $\mathcal F$.
There then exist $C'_B,t_1,\rho_1>0$ such that, after centering at the nonlinear
conditional minimizer and restricting to a common convex fiber ball,
\begin{equation}\label{eq:cross}
 t\sup_{w\in D}
 \|\nabla V_y(w)-\nabla V_{y'}(w)\|
 \le \chi\,d_Y(y,y'),
 \qquad \chi\le C'_B\rho.
\end{equation}
for $0<t\le t_1$ and a sufficiently small tube of radius $\rho\le\rho_1$.
\end{lemma}
\begin{proof}
Choose the horizontal derivative of $\Phi$ at the origin to be the
$L$-orthogonal conditional-mean map $M$ from
Theorem~\ref{thm:gaussian}; choose its vertical derivative to identify
$\R^{357}$ with $\mathcal F$.  The mixed base--fiber Hessian of the leading
quadratic action then vanishes.  Centering at the nonlinear minimizer also
makes the fiber gradient vanish at $w=0$.  After subtracting the quadratic
heat-kernel term, the map
\[
 (t,y,w)\longmapsto t\nabla_wV_y(w)
\]
has uniformly bounded derivatives on a smaller local tube by the same
parametrix and BCH argument as Lemma~\ref{lem:hk}; exponentially small image
terms are absorbed by decreasing $t_1$.  Its base derivative at the origin is
zero.  The mean-value theorem therefore bounds its base derivative by
$C'_B\rho$, proving \eqref{eq:cross}.
\end{proof}

The factor $\rho$ records that the leading Gaussian cross-response has been
removed; what remains begins with nonlinear commutator and coordinate terms.

\begin{corollary}[Local centered-fiber contraction, \proved\ (local)]
If \eqref{eq:hessian-remainder} and \eqref{eq:cross} hold with
\[
 K_{\mathrm{fiber}}:=\frac{\chi}{2-\eta}<1,
\]
then
\[
 W_2(\widetilde\nu_y,\widetilde\nu_{y'})
 \le K_{\mathrm{fiber}}d_Y(y,y').
\]
For the fixed block, sufficiently small $\rho$ and $t$ make this inequality
true with existential constants.
\end{corollary}
\begin{proof}
Use $\alpha=(2-\eta)/t$ from Theorem~\ref{thm:nonlinear} and
$L=\chi/t$ from Lemma~\ref{lem:cross} in Lemma~\ref{lem:transport}.  Since
$\eta\to0$ and $\chi\to0$ as $(\rho,t)\to(0,0)$ inside the local tube, shrink
the tube until $\chi/(2-\eta)<1$.
\end{proof}

The common-domain restriction is deliberate.  The probability that leaves
this good tube belongs to the defect channel; it cannot be discarded.

\section{Executable calibration and hostile tests}

The accompanying implementation constructs $d_0,d_1,C$, uses exact rational
elimination for the retained rank, and computes orthogonal quotient and fiber
bases numerically.  Its output is:

\begin{center}
\begin{tabular}{@{}lr@{}}
\toprule
Quantity & Value\\
\midrule
scalar gauge rank & $80$\\
scalar physical dimension & $136$\\
displayed / independent coarse fluxes & $24/17$\\
scalar fiber dimension & $119$\\
$\su$ physical / base / fiber dimensions & $408/51/357$\\
physical eigenvalue range & $[2,12]$\\
fiber eigenvalue range & $[3.438447187,12]$\\
minimum nonzero coarse singular value & $2.828427125$\\
full conditional-mean response norm & $0.393174006$\\
centered Gaussian response & $0$\\
\bottomrule
\end{tabular}
\end{center}

Eleven automated tests include four hostile checks:
\begin{enumerate}
\item without quotienting, the quadratic Hessian has exactly $80$ scalar
      gauge zero modes;
\item the $24$ displayed retained fluxes are not falsely treated as
      independent;
\item invalid geometries and base dimensions fail closed;
\item a nonlinear remainder that exhausts the certified margin is rejected.
\end{enumerate}

The numerical eigenvalue and response norm calibrate the theorem; only the
analytic lower bound $2$, exact ranks, and Gaussian identities enter proved
claims.

\section{RG certificate and large-field ledger}

A complete block morphism would transport
\[
 (\mu_k,S_k,\mathcal O_k,\ell_k)
 \longmapsto
 (\mu_{k+1},S_{k+1},\mathcal O_{k+1},b\ell_k)
\]
together with the retained map, conditional fiber, local interaction norm,
running coupling, reconstruction error, and positivity certificate.  The
present package supplies only the finite local tangent and good-tube entries.

\begin{longtable}{@{}p{0.08\textwidth}p{0.62\textwidth}p{0.18\textwidth}@{}}
\toprule
ID & Required result & Status\\
\midrule
\endhead
F1 & Replace the local Coulomb/constant-rank chart by an exact gauge-covariant
nonlinear disintegration, with orbit-atlas overlaps and stabilizers. & \open\\
F2 & Derive explicit $C_B,C'_B,\rho_0,t_0$ and prove stability for the full
generated effective-action class under one RG step. & \open\\
F3 & Bound the complement of the good tube by a dependent polymer norm,
including cut-locus configurations, reducible orbits, center flux, global
holonomy, and topological sectors. & \open\\
F4 & Iterate through $k(a)\sim\log_b(\ell_*/a)\to\infty$ ultraviolet steps
with summable errors and controlled observable support. & \open\\
F5 & Prove that the retained base crosses into a massive basin at a finite
physical scale. & \open\\
F6 & Construct nontrivial regulator-independent continuum Schwinger functions,
verify the OS hypotheses and observable density, and transfer decay to the
Hamiltonian spectrum. & \open\\
\bottomrule
\end{longtable}

Reflection positivity need not be manifest in every intermediate coordinate
system if the RG integration is exact and the original regulated observables
are transported without error; otherwise it requires a quantitative recovery
certificate.  No such multiscale certificate is supplied here.

\section{What has and has not advanced}

The result establishes that the retained/fiber distinction survives the first
genuinely four-dimensional nonabelian-compatible test.  Gauge directions are
removed exactly at tangent level, retained curvature data have the correct
rank, the fiber has a certified coercivity margin, and nonlinear heat-kernel
geometry preserves that margin locally.  The centered transport coefficient
is therefore a mathematically viable local target rather than a purely
ontological proposal.

The result does not generate infrared mass.  Fiber contraction controls the
degrees of freedom discarded by blocking; the retained base can still carry
massless long-wavelength modes.  The central Yang--Mills problem is now located
more sharply in F4--F5: control a diverging RG trajectory and prove a
nonperturbative crossover of the base into a massive regime.  The local theorem
is an input to that task, not its conclusion.

\section{Immediate next theorem}

The next bounded theorem should attach the certificate to a Balaban-type
effective-action norm.  For one exact RG step, prove
\[
 t_k\|\nabla_F^2(S_k-Q_k)\|\le\eta_k<2,
 \qquad
 t_k\|\nabla_FV_{k,y}-\nabla_FV_{k,y'}\|
 \le\chi_k d_Y(y,y'),
\]
with
\[
 \frac{\chi_k}{2-\eta_k}<1,
\]
and a gauge-invariant defect activity satisfying a declared polymer norm.
Every constant must be expressed in variables propagated by the RG map.  A
failure of this inequality would identify whether the obstruction lies in the
chosen retained data, the orbit chart, or the generated interactions.

\begin{thebibliography}{9}
\bibitem{JaffeWitten}
A.~Jaffe and E.~Witten,
\emph{Quantum Yang--Mills Theory}, Clay Mathematics Institute Millennium
Problem description,
\url{https://www.claymath.org/wp-content/uploads/2022/06/yangmills.pdf}.

\bibitem{MenottiOnofri}
P.~Menotti and E.~Onofri,
``The action of $SU(N)$ lattice gauge theory in terms of the heat kernel on the
group manifold,'' \emph{Nuclear Physics B} \textbf{190} (1981), 288--300,
doi:10.1016/0550-3213(81)90560-5.

\bibitem{BalabanBackground}
T.~Balaban,
``The variational problem and background fields in renormalization group
method for lattice gauge theories,''
\emph{Communications in Mathematical Physics} \textbf{102} (1985), 277--309,
doi:10.1007/BF01229381.

\bibitem{BalabanPropagators}
T.~Balaban,
``Propagators and renormalization transformations for lattice gauge theories
I,'' \emph{Communications in Mathematical Physics} \textbf{95} (1984), 17--40,
doi:10.1007/BF01215753.

\bibitem{OttoVillani}
F.~Otto and C.~Villani,
``Generalization of an inequality by Talagrand and links with the logarithmic
Sobolev inequality,'' \emph{Journal of Functional Analysis} \textbf{173}
(2000), 361--400, doi:10.1006/jfan.1999.3557.
\end{thebibliography}

\end{document}
